\documentclass[UTF8]{ctexart}

\usepackage{amsmath,amssymb}
\usepackage[a4paper,margin=2.4cm]{geometry}

\title{反比例函数、极坐标与等轴双曲线}
\author{}
\date{}

\begin{document}

\maketitle

\section{从反比例函数到曲线方程}

给定反比例函数
\[
  y=\frac{k}{x},\qquad k\ne 0,
\]
把分母移到等式左边，可以写成直角坐标方程
\[
  xy=k.
\]
这条曲线的两条渐近线是
\[
  x=0,\qquad y=0,
\]
也就是坐标轴本身。两条坐标轴互相垂直，所以从渐近线的角度看，
这条曲线已经具有等轴双曲线的特征。下面用极坐标和旋转变换把这一点写成标准形式。

\section{化为极坐标方程}

在极坐标中，
\[
  x=r\cos\theta,\qquad y=r\sin\theta.
\]
代入直角坐标方程 \(xy=k\)，得到
\[
  (r\cos\theta)(r\sin\theta)=k.
\]
因此
\[
  r^2\sin\theta\cos\theta=k.
\]
利用二倍角公式
\[
  \sin 2\theta=2\sin\theta\cos\theta,
\]
可得
\[
  \frac{r^2}{2}\sin 2\theta=k,
\]
即
\[
  r^2\sin 2\theta=2k.
\]
如果写成 \(r^2\) 的形式，就是
\[
  r^2=\frac{2k}{\sin 2\theta}.
\]
这就是反比例函数 \(y=k/x\) 对应的极坐标方程。

\section{旋转极轴}

现在把极轴逆时针旋转 \(45^\circ\)，即旋转角
\[
  \alpha=\frac{\pi}{4}.
\]
同一个点到原点的距离 \(r\) 不变，只是角度的表示发生变化。若新极坐标角记为
\[
  \varphi=\theta-\frac{\pi}{4},
\]
则
\[
  \theta=\varphi+\frac{\pi}{4}.
\]
把它代入极坐标方程
\[
  r^2\sin 2\theta=2k,
\]
得到
\[
  r^2\sin\left(2\varphi+\frac{\pi}{2}\right)=2k.
\]
因为
\[
  \sin\left(2\varphi+\frac{\pi}{2}\right)=\cos 2\varphi,
\]
所以旋转后的方程变为
\[
  r^2\cos 2\varphi=2k.
\]

\section{得到标准等轴双曲线}

在旋转后的坐标系中，记
\[
  X=r\cos\varphi,\qquad Y=r\sin\varphi.
\]
于是
\[
  r^2\cos 2\varphi
  =r^2(\cos^2\varphi-\sin^2\varphi)
  =X^2-Y^2.
\]
因此旋转后的方程
\[
  r^2\cos 2\varphi=2k
\]
等价于
\[
  X^2-Y^2=2k.
\]

当 \(k>0\) 时，可以写成
\[
  \frac{X^2}{2k}-\frac{Y^2}{2k}=1.
\]
它的实半轴和虚半轴相等，均为
\[
  \sqrt{2k}.
\]
所以它是等轴双曲线。

当 \(k<0\) 时，方程可以写成
\[
  Y^2-X^2=-2k,
\]
也就是
\[
  \frac{Y^2}{-2k}-\frac{X^2}{-2k}=1.
\]
此时两条半轴仍然相等，均为
\[
  \sqrt{-2k}.
\]
所以它仍然是等轴双曲线，只是开口方向与 \(k>0\) 的情形相反。

\section{结论}

反比例函数
\[
  y=\frac{k}{x}\qquad (k\ne 0)
\]
对应的直角坐标方程是
\[
  xy=k.
\]
把它化为极坐标方程，得到
\[
  r^2\sin 2\theta=2k.
\]
再将极轴旋转 \(45^\circ\)，令
\[
  \varphi=\theta-\frac{\pi}{4},
\]
方程变为
\[
  r^2\cos 2\varphi=2k.
\]
回到旋转后的直角坐标系，就得到
\[
  X^2-Y^2=2k.
\]
这正是实半轴和虚半轴相等的标准双曲线形式，因此反比例函数的图像本质上是一条等轴双曲线。

\end{document}
